% A non-effective finite equivalence relation on a finite scheme over a field — standalone article for Hexagon (problem of the Buffalo lecture of 11 July 1973; FGA 3, exp. 212, §2).
% pdfLaTeX. Compile: latexmk -pdf article.tex (or tectonic -X compile article.tex).
\documentclass[11pt,a4paper]{amsart}
\usepackage[T1]{fontenc}
\usepackage[utf8]{inputenc}
\usepackage{lmodern,microtype,amssymb,mathtools,booktabs}
\usepackage[margin=2.7cm]{geometry}
\usepackage[hidelinks]{hyperref}

\newtheorem{theorem}{Theorem}
\newtheorem{proposition}[theorem]{Proposition}
\newtheorem{lemma}[theorem]{Lemma}
\theoremstyle{remark}
\newtheorem{remark}[theorem]{Remark}

\newcommand{\A}{\mathbb A}
\newcommand{\F}{\mathbb F}
\newcommand{\Q}{\mathbb Q}

\title[A non-effective finite equivalence relation]{A non-effective finite equivalence relation\\ on a finite scheme over a field,\\ after a question of Grothendieck}
\author{Michel Hua}
\address{Commutator, Paris, France}
\email{michel@commutator.io}
\thanks{ORCID 0009-0002-8304-0683.}
\date{7 November 2026}
\subjclass[2020]{14L30; 14A15, 13A50, 13D03}
\keywords{finite equivalence relations, effectivity, quotients, descent, Artin rings, Grothendieck}

\begin{document}
\begin{abstract}
In 1961 Grothendieck asked whether a closed equivalence relation on a finite scheme over a field, with finite but not necessarily flat projections, is always effective, and in a lecture at Buffalo in 1973 he still held the question open. We show that the answer is no over every field. Restricting Raicu's non-effective relation on the affine plane to a finite subscheme of length $27$ cut out by powers of invariant functions gives a closed equivalence relation on the spectrum of a local Artin algebra with residue field $k$ which is not effective: its ideal is strictly larger, in degrees $5$ and $6$, than the ideal of $X \times_Y X$. The comparison is a finite computation in linear algebra; we carry it out over $\Q$, record the primes that divide a pivot, and check those characteristics separately, so that it holds over every field.
\end{abstract}
\maketitle

\section{The question}

Let $k$ be a field, $X = \operatorname{Spec} A$ a finite $k$-scheme and $R \subset X \times_k X$ a closed subscheme whose points with values in every $k$-algebra form an equivalence relation. Its projections $R \rightrightarrows X$ are finite, but they need not be flat. The ring of invariants
\[
B = A^R = \{\, a \in A : a \otimes 1 - 1 \otimes a \in I_R \,\},
\]
with $I_R \subset A \otimes_k A$ the ideal of $R$, is a finite $k$-algebra, and $Y = \operatorname{Spec} B$ is the quotient of $X$ by $R$ in the category of schemes. The relation is \emph{effective} when the morphism $R \to X \times_Y X$ is an isomorphism.

In the third of his talks on the construction of schemes, after the flat case, Grothendieck wrote that without flatness he did not know whether a finite equivalence relation is effective. He reduced the question to the spectrum of a local Artin ring with algebraically closed residue field, and said that it was open even over a field \cite[\S2, p.~212-05]{fga3}. Twelve years later, at the end of a lecture of his algebraic geometry course at the State University of New York at Buffalo, on 11 July 1973, a listener asked for open problems and he gave this one: he said it was not known, that he expected the answer to be negative, and that it was very elementary \cite{buffalo}. The flat case is SGA~3, exp.~V \cite{sga3}.

Without flatness, two kinds of negative results are known. Venken showed that descent of flat modules along a finite morphism of local Artin rings need not be effective \cite{venken}, and Philippe studied descent along unramified morphisms \cite{philippe}; these concern modules, not the relation itself. In the appendix to Kollár's paper on quotients by finite equivalence relations, Raicu gave finite equivalence relations on $\A^2$ which are not effective \cite[\S50]{kollar}, while affine toric equivalence relations are effective \cite{raicu2010}. We have not found the case of a finite scheme over a field treated in print, in either direction.

\begin{theorem}\label{thm:main}
Let $k$ be any field. There is a local Artin $k$-algebra $A$ of dimension $27$, with residue field $k$, and a closed equivalence relation $R$ on $\operatorname{Spec} A$, with finite projections, which is not effective.
\end{theorem}

Over a field the answer is therefore negative, as Grothendieck expected, and the more general question, for $X$ finite over a noetherian base, is answered by the same example.

\section{Effectivity as linear algebra}

With the notation of the question, $X \times_Y X = \operatorname{Spec}(A \otimes_B A)$ and $A \otimes_B A = (A \otimes_k A)/J_B$, where $J_B$ is the ideal generated by the $b \otimes 1 - 1 \otimes b$, $b \in B$. By definition of $B$, $J_B \subseteq I_R$, and $R \to X \times_Y X$ is the closed immersion defined by $I_R / J_B$. So $R$ is effective if and only if $I_R = J_B$, a comparison of two subspaces of the finite-dimensional space $A \otimes_k A$.

When $A$ and $I_R$ are graded, $B$ is a graded subalgebra and both ideals are graded, so the comparison can be made degree by degree. A homogeneous invariant of degree $e$ contributes to $J_B$ only in degrees $\ge e$; to compare the two ideals in degree $d$ it suffices to know the invariants of degree at most $d$.

\section{Raicu's relation and its restriction}

Write $x = (x_1, x_2)$ and $y = (y_1, y_2)$ for the coordinates of $\A^2 \times \A^2$. Raicu's relation \cite[\S50]{kollar} is the closed subscheme $R \subset \A^2 \times \A^2$ with ideal
\[
I = \bigl(x_1^2 - y_1^2,\ x_1x_2 - x_2^2 - y_1y_2 + y_2^2,\ x_2^3 - y_2^3,\ (x_1y_2 - x_2y_1)\,y_2^3\bigr) \subset k[x_1, x_2, y_1, y_2].
\]
It is a finite equivalence relation on $\A^2$, defined over $\mathbb Z$, and it is not effective after base change to any field; the class of $(x_1y_2 - x_2y_1)\,y_2^3$ is the obstruction. The functions
\[
f_1 = x_1^2, \qquad f_2 = x_1x_2 - x_2^2, \qquad f_3 = x_2^3
\]
are invariant, since $f_i(x) - f_i(y)$ is one of the first three generators of $I$. Let
\[
Z = V(f_1^3, f_2^3, f_3^2) \subset \A^2, \qquad A = k[x_1, x_2]/(f_1^3, f_2^3, f_3^2).
\]
Since $f_1^3 = x_1^6$ and $f_3^2 = x_2^6$, $Z$ is supported at the origin, and $A$ is a local Artin algebra with residue field $k$; a count of monomials shows that $\dim_k A = 27$. Let $R_Z = R \cap (Z \times Z)$, the closed subscheme of $Z \times Z$ with ideal $I_{R_Z}$, the image of $I$ in $A \otimes_k A$.

\begin{lemma}
$R_Z$ is a closed equivalence relation on $Z$ with finite projections.
\end{lemma}

\begin{proof}
For every $k$-algebra $T$, $R_Z(T) = R(T) \cap (Z(T) \times Z(T))$ is the restriction of an equivalence relation on $\A^2(T)$ to the subset $Z(T)$, hence an equivalence relation. $R_Z$ is closed in $Z \times Z$, which is finite over $Z$ by either projection.
\end{proof}

The restriction to $Z$ is what lets a question about $\A^2$ be asked about a finite scheme. It is not formal that non-effectivity survives it: the ring of invariants of $R_Z$ is computed on $A$, not on $k[x_1, x_2]$, and it can be larger than the image of the invariants on $\A^2$, so that $J_B$ can grow. The next section shows that it does not grow enough.

\section{The computation}

All the ideals above are homogeneous for the standard grading. In $k[x, y] = k[x_1, x_2, y_1, y_2]$ put $K = (f_1^3, f_2^3, f_3^2)(x) + (f_1^3, f_2^3, f_3^2)(y)$, so that $A \otimes_k A = k[x, y]/K$. The ideal of $R_Z$ lifts to $I + K$, and $J_B$ lifts to $J + K$, where $J$ is generated by the $b(x) - b(y)$ for $b$ running over homogeneous lifts of a basis of $B$.

\begin{proposition}\label{prop:dims}
Over $\Q$ and over every finite prime field, the invariants of $R_Z$ of degree at most $7$ give $22$ independent differences $b(x) - b(y)$ modulo $I + K$, and in the degrees $d \le 14$ the two ideals differ exactly in
\[
\dim_k (I + K)_5 = 44 > 43 = \dim_k (J + K)_5, \qquad \dim_k (I + K)_6 = 77 > 75 = \dim_k (J + K)_6.
\]
\end{proposition}

By the remark at the end of Section~2, invariants of degree at least $8$ do not change $J + K$ in degrees $5$ and $6$, so the strict inequality in degree $5$ proves that $R_Z$ is not effective. A witness is the generator $(x_1y_2 - x_2y_1)\,y_2^3$ of $I$, Raicu's obstruction: it lies in $(I + K)_5$ and not in $(J + K)_5$, as we checked over $\Q$, $\F_2$, $\F_3$ and $\F_5$.

The proposition is a finite computation: the degree-$d$ parts are spanned by monomial multiples of the generators, the invariants of degree $d$ are the kernel of a linear map on the forms of degree $d$ in $x$, and the dimensions are ranks. The script \texttt{full.py} of the directory \texttt{computations/effective/raicu} of the repository \cite{repo} performs it by Gaussian elimination over $\Q$ or over $\F_p$.

\begin{proof}[Proof that it holds over every field]
A field $k$ contains $\Q$ or some $\F_p$, and all the data are defined over the prime field: ranks and kernels of matrices with entries in the prime field do not change under extension of scalars. It remains to cover every prime field.

Over $\Q$ the elimination divides only by pivots. Each coefficient it produces is an integer combination of the input coefficients, all integers, divided by a product of pivots. If a prime $p$ divides neither the numerator nor the denominator of any pivot, the same sequence of operations is defined over $\F_p$, with the same pivot positions, and gives the same ranks and the same kernels. The script \texttt{certify.py} reruns the computation over $\Q$ and records the prime factors of all pivots; they are $2, 3, 5, 7, 11, 13, 41, 61, 67$ and $661$. For every other prime the proposition over $\F_p$ follows from the one over $\Q$. For these ten primes, and for all primes up to $101$ and for $1009$ and $10007$ as a cross-check, \texttt{full.py} was run directly over $\F_p$, with the same result.
\end{proof}

\begin{remark}
Why $Z$ and not something smaller. Before this example, a random search over $48$ small monomial algebras over $\F_2$ and $\F_3$ found no non-effective relation. Raicu's theorem that affine toric equivalence relations are effective \cite{raicu2010} suggests why: monomial relations are close to the toric case. We did not look for the smallest finite subscheme on which the obstruction survives.
\end{remark}

\begin{remark}
Grothendieck reduced the question to an algebraically closed residue field; the example has residue field $k$, which may be chosen algebraically closed. It is also graded, which the question does not require.
\end{remark}

\subsection*{What is not done}
The computation is not formalised: a Lean proof would need either a certified Gaussian elimination over $\F_3$ and $\Q$ or a Gröbner-basis certificate, and the pivot argument above. The example has not been checked by a specialist of descent.

\subsection*{Use of language models}
The question was located in the recording of the Buffalo lecture and in \cite{fga3} with language models (Claude Opus~5.5, Anthropic), which also found Raicu's example, proposed the restriction to $Z$, wrote the scripts and a first draft of this text, under the author's direction. The proposition rests on the scripts, which anyone can rerun.

\begin{thebibliography}{9}
\bibitem{buffalo} A.~Grothendieck, course on algebraic groups (with algebraic geometry), State University of New York at Buffalo, lecture of 11 July 1973, recording at 6:59:50--7:00:54 of \url{https://www.youtube.com/watch?v=3ULFIl_2Wis}.
\bibitem{fga3} A.~Grothendieck, Techniques de construction et théorèmes d'existence en géométrie algébrique~III : préschémas quotients, Séminaire Bourbaki~6 (1960/61), exp.~212, 99--118, \url{http://www.numdam.org/item/SB_1960-1961__6__99_0}.
\bibitem{kollar} J.~Kollár, Quotients by finite equivalence relations, with an appendix by C.~Raicu, in \emph{Current Developments in Algebraic Geometry}, MSRI Publ.~59, Cambridge University Press, 2012, 227--256, \href{https://arxiv.org/abs/0812.3608}{arXiv:0812.3608}.
\bibitem{philippe} A.~Philippe, Morphisme net d'anneaux, et descente, \emph{Bull.\ Sci.\ Math.~(2)} 97 (1973), 57--64.
\bibitem{raicu2010} C.~Raicu, Affine toric equivalence relations are effective, \emph{Proc.\ Amer.\ Math.\ Soc.} 138 (2010), 3835--3847, \href{https://doi.org/10.1090/S0002-9939-2010-10416-1}{doi:10.1090/S0002-9939-2010-10416-1}.
\bibitem{repo} Commutator-IO, grothendieck-project, \url{https://github.com/Commutator-IO/grothendieck-project}, directory \texttt{computations/effective/raicu}.
\bibitem{sga3} M.~Demazure, A.~Grothendieck (eds.), \emph{Schémas en groupes} (SGA~3), Lecture Notes in Math.~151, 152, 153, Springer, 1970.
\bibitem{venken} J.~Venken, Non effectivité de la descente de modules plats par un morphisme fini d'anneaux locaux artiniens, \emph{C.\ R.\ Acad.\ Sci.\ Paris Sér.~A--B} 272 (1971), A1553--A1554.
\end{thebibliography}
\end{document}
